\documentclass[11pt]{article}
\usepackage[left=1.5cm,top=1.5cm,right=1.5cm, bottom=1.5cm,letterpaper, includeheadfoot]{geometry}
\usepackage{hyperref}
\hypersetup{
	colorlinks,
	citecolor=black,
	filecolor=black,
	linkcolor=black,
	urlcolor=black
}

\usepackage{amssymb, amsmath, amsthm}
\usepackage{graphicx}
\usepackage{lmodern,url}


\usepackage[english]{babel}
\usepackage{fancyhdr}
\usepackage{amsmath}
\usepackage{lipsum}
\usepackage{framed}
\usepackage{tikz}
\usepackage{bbm}
\usepackage{graphicx,todo}
\graphicspath{{images/}}
\usetikzlibrary{matrix}
\usepackage{multicol}






\usepackage[breakable]{tcolorbox}
\usepackage{parskip} % Stop auto-indenting (to mimic markdown behaviour)


% Basic figure setup, for now with no caption control since it's done
% automatically by Pandoc (which extracts ![](path) syntax from Markdown).

% Maintain compatibility with old templates. Remove in nbconvert 6.0
\let\Oldincludegraphics\includegraphics
% Ensure that by default, figures have no caption (until we provide a
% proper Figure object with a Caption API and a way to capture that
% in the conversion process - todo).
\usepackage{caption}
% \DeclareCaptionFormat{nocaption}{}
% \captionsetup{format=nocaption,aboveskip=0pt,belowskip=0pt}

\usepackage{float}
\floatplacement{figure}{H} % forces figures to be placed at the correct location
\usepackage{xcolor} % Allow colors to be defined
\usepackage{enumerate} % Needed for markdown enumerations to work
\usepackage{geometry} % Used to adjust the document margins

\usepackage{amssymb} % Equations
\usepackage{textcomp} % defines textquotesingle
% Hack from http://tex.stackexchange.com/a/47451/13684:
\AtBeginDocument{%
    \def\PYZsq{\textquotesingle}% Upright quotes in Pygmentized code
}
\usepackage{upquote} % Upright quotes for verbatim code
\usepackage{eurosym} % defines \euro

\usepackage{iftex}
\ifPDFTeX
    \usepackage[T1]{fontenc}
    \IfFileExists{alphabeta.sty}{
          \usepackage{alphabeta}
      }{
          \usepackage[mathletters]{ucs}
          \usepackage[utf8x]{inputenc}
      }
\else
    \usepackage{fontspec}
    \usepackage{unicode-math}
\fi

\usepackage{fancyvrb} % verbatim replacement that allows latex
\usepackage{grffile} % extends the file name processing of package graphics 
                     % to support a larger range
\makeatletter % fix for old versions of grffile with XeLaTeX
\@ifpackagelater{grffile}{2019/11/01}
{
  % Do nothing on new versions
}
{
  \def\Gread@@xetex#1{%
    \IfFileExists{"\Gin@base".bb}%
    {\Gread@eps{\Gin@base.bb}}%
    {\Gread@@xetex@aux#1}%
  }
}
\makeatother
\usepackage[Export]{adjustbox} % Used to constrain images to a maximum size
\adjustboxset{max size={0.9\linewidth}{0.9\paperheight}}

% The hyperref package gives us a pdf with properly built
% internal navigation ('pdf bookmarks' for the table of contents,
% internal cross-reference links, web links for URLs, etc.)
\usepackage{hyperref}
% The default LaTeX title has an obnoxious amount of whitespace. By default,
% titling removes some of it. It also provides customization options.
\usepackage{titling}
\usepackage{longtable} % longtable support required by pandoc >1.10
\usepackage{booktabs}  % table support for pandoc > 1.12.2
\usepackage{array}     % table support for pandoc >= 2.11.3
\usepackage{calc}      % table minipage width calculation for pandoc >= 2.11.1
\usepackage[inline]{enumitem} % IRkernel/repr support (it uses the enumerate* environment)
\usepackage[normalem]{ulem} % ulem is needed to support strikethroughs (\sout)
                            % normalem makes italics be italics, not underlines
\usepackage{mathrsfs}



% Colors for the hyperref package
\definecolor{urlcolor}{rgb}{0,.145,.698}
\definecolor{linkcolor}{rgb}{.71,0.21,0.01}
\definecolor{citecolor}{rgb}{.12,.54,.11}

% ANSI colors
\definecolor{ansi-black}{HTML}{3E424D}
\definecolor{ansi-black-intense}{HTML}{282C36}
\definecolor{ansi-red}{HTML}{E75C58}
\definecolor{ansi-red-intense}{HTML}{B22B31}
\definecolor{ansi-green}{HTML}{00A250}
\definecolor{ansi-green-intense}{HTML}{007427}
\definecolor{ansi-yellow}{HTML}{DDB62B}
\definecolor{ansi-yellow-intense}{HTML}{B27D12}
\definecolor{ansi-blue}{HTML}{208FFB}
\definecolor{ansi-blue-intense}{HTML}{0065CA}
\definecolor{ansi-magenta}{HTML}{D160C4}
\definecolor{ansi-magenta-intense}{HTML}{A03196}
\definecolor{ansi-cyan}{HTML}{60C6C8}
\definecolor{ansi-cyan-intense}{HTML}{258F8F}
\definecolor{ansi-white}{HTML}{C5C1B4}
\definecolor{ansi-white-intense}{HTML}{A1A6B2}
\definecolor{ansi-default-inverse-fg}{HTML}{FFFFFF}
\definecolor{ansi-default-inverse-bg}{HTML}{000000}

% common color for the border for error outputs.
\definecolor{outerrorbackground}{HTML}{FFDFDF}

% commands and environments needed by pandoc snippets
% extracted from the output of `pandoc -s`
\providecommand{\tightlist}{%
  \setlength{\itemsep}{0pt}\setlength{\parskip}{0pt}}
\DefineVerbatimEnvironment{Highlighting}{Verbatim}{commandchars=\\\{\}}
% Add ',fontsize=\small' for more characters per line
\newenvironment{Shaded}{}{}
\newcommand{\KeywordTok}[1]{\textcolor[rgb]{0.00,0.44,0.13}{\textbf{{#1}}}}
\newcommand{\DataTypeTok}[1]{\textcolor[rgb]{0.56,0.13,0.00}{{#1}}}
\newcommand{\DecValTok}[1]{\textcolor[rgb]{0.25,0.63,0.44}{{#1}}}
\newcommand{\BaseNTok}[1]{\textcolor[rgb]{0.25,0.63,0.44}{{#1}}}
\newcommand{\FloatTok}[1]{\textcolor[rgb]{0.25,0.63,0.44}{{#1}}}
\newcommand{\CharTok}[1]{\textcolor[rgb]{0.25,0.44,0.63}{{#1}}}
\newcommand{\StringTok}[1]{\textcolor[rgb]{0.25,0.44,0.63}{{#1}}}
\newcommand{\CommentTok}[1]{\textcolor[rgb]{0.38,0.63,0.69}{\textit{{#1}}}}
\newcommand{\OtherTok}[1]{\textcolor[rgb]{0.00,0.44,0.13}{{#1}}}
\newcommand{\AlertTok}[1]{\textcolor[rgb]{1.00,0.00,0.00}{\textbf{{#1}}}}
\newcommand{\FunctionTok}[1]{\textcolor[rgb]{0.02,0.16,0.49}{{#1}}}
\newcommand{\RegionMarkerTok}[1]{{#1}}
\newcommand{\ErrorTok}[1]{\textcolor[rgb]{1.00,0.00,0.00}{\textbf{{#1}}}}
\newcommand{\NormalTok}[1]{{#1}}

% Additional commands for more recent versions of Pandoc
\newcommand{\ConstantTok}[1]{\textcolor[rgb]{0.53,0.00,0.00}{{#1}}}
\newcommand{\SpecialCharTok}[1]{\textcolor[rgb]{0.25,0.44,0.63}{{#1}}}
\newcommand{\VerbatimStringTok}[1]{\textcolor[rgb]{0.25,0.44,0.63}{{#1}}}
\newcommand{\SpecialStringTok}[1]{\textcolor[rgb]{0.73,0.40,0.53}{{#1}}}
\newcommand{\ImportTok}[1]{{#1}}
\newcommand{\DocumentationTok}[1]{\textcolor[rgb]{0.73,0.13,0.13}{\textit{{#1}}}}
\newcommand{\AnnotationTok}[1]{\textcolor[rgb]{0.38,0.63,0.69}{\textbf{\textit{{#1}}}}}
\newcommand{\CommentVarTok}[1]{\textcolor[rgb]{0.38,0.63,0.69}{\textbf{\textit{{#1}}}}}
\newcommand{\VariableTok}[1]{\textcolor[rgb]{0.10,0.09,0.49}{{#1}}}
\newcommand{\ControlFlowTok}[1]{\textcolor[rgb]{0.00,0.44,0.13}{\textbf{{#1}}}}
\newcommand{\OperatorTok}[1]{\textcolor[rgb]{0.40,0.40,0.40}{{#1}}}
\newcommand{\BuiltInTok}[1]{{#1}}
\newcommand{\ExtensionTok}[1]{{#1}}
\newcommand{\PreprocessorTok}[1]{\textcolor[rgb]{0.74,0.48,0.00}{{#1}}}
\newcommand{\AttributeTok}[1]{\textcolor[rgb]{0.49,0.56,0.16}{{#1}}}
\newcommand{\InformationTok}[1]{\textcolor[rgb]{0.38,0.63,0.69}{\textbf{\textit{{#1}}}}}
\newcommand{\WarningTok}[1]{\textcolor[rgb]{0.38,0.63,0.69}{\textbf{\textit{{#1}}}}}

\newcommand{\wrong}[1] {\textcolor{red}{[#1]}}
% Define a nice break command that doesn't care if a line doesn't already
% exist.
\def\br{\hspace*{\fill} \\* }
% Math Jax compatibility definitions
\def\gt{>}
\def\lt{<}
\let\Oldtex\TeX
\let\Oldlatex\LaTeX
\renewcommand{\TeX}{\textrm{\Oldtex}}
\renewcommand{\LaTeX}{\textrm{\Oldlatex}}
% Document parameters
% Document title
\title{Numerical Alg}





% Pygments definitions
\makeatletter
\def\PY@reset{\let\PY@it=\relax \let\PY@bf=\relax%
\let\PY@ul=\relax \let\PY@tc=\relax%
\let\PY@bc=\relax \let\PY@ff=\relax}
\def\PY@tok#1{\csname PY@tok@#1\endcsname}
\def\PY@toks#1+{\ifx\relax#1\empty\else%
\PY@tok{#1}\expandafter\PY@toks\fi}
\def\PY@do#1{\PY@bc{\PY@tc{\PY@ul{%
\PY@it{\PY@bf{\PY@ff{#1}}}}}}}
\def\PY#1#2{\PY@reset\PY@toks#1+\relax+\PY@do{#2}}

\@namedef{PY@tok@w}{\def\PY@tc##1{\textcolor[rgb]{0.73,0.73,0.73}{##1}}}
\@namedef{PY@tok@c}{\let\PY@it=\textit\def\PY@tc##1{\textcolor[rgb]{0.24,0.48,0.48}{##1}}}
\@namedef{PY@tok@cp}{\def\PY@tc##1{\textcolor[rgb]{0.61,0.40,0.00}{##1}}}
\@namedef{PY@tok@k}{\let\PY@bf=\textbf\def\PY@tc##1{\textcolor[rgb]{0.00,0.50,0.00}{##1}}}
\@namedef{PY@tok@kp}{\def\PY@tc##1{\textcolor[rgb]{0.00,0.50,0.00}{##1}}}
\@namedef{PY@tok@kt}{\def\PY@tc##1{\textcolor[rgb]{0.69,0.00,0.25}{##1}}}
\@namedef{PY@tok@o}{\def\PY@tc##1{\textcolor[rgb]{0.40,0.40,0.40}{##1}}}
\@namedef{PY@tok@ow}{\let\PY@bf=\textbf\def\PY@tc##1{\textcolor[rgb]{0.67,0.13,1.00}{##1}}}
\@namedef{PY@tok@nb}{\def\PY@tc##1{\textcolor[rgb]{0.00,0.50,0.00}{##1}}}
\@namedef{PY@tok@nf}{\def\PY@tc##1{\textcolor[rgb]{0.00,0.00,1.00}{##1}}}
\@namedef{PY@tok@nc}{\let\PY@bf=\textbf\def\PY@tc##1{\textcolor[rgb]{0.00,0.00,1.00}{##1}}}
\@namedef{PY@tok@nn}{\let\PY@bf=\textbf\def\PY@tc##1{\textcolor[rgb]{0.00,0.00,1.00}{##1}}}
\@namedef{PY@tok@ne}{\let\PY@bf=\textbf\def\PY@tc##1{\textcolor[rgb]{0.80,0.25,0.22}{##1}}}
\@namedef{PY@tok@nv}{\def\PY@tc##1{\textcolor[rgb]{0.10,0.09,0.49}{##1}}}
\@namedef{PY@tok@no}{\def\PY@tc##1{\textcolor[rgb]{0.53,0.00,0.00}{##1}}}
\@namedef{PY@tok@nl}{\def\PY@tc##1{\textcolor[rgb]{0.46,0.46,0.00}{##1}}}
\@namedef{PY@tok@ni}{\let\PY@bf=\textbf\def\PY@tc##1{\textcolor[rgb]{0.44,0.44,0.44}{##1}}}
\@namedef{PY@tok@na}{\def\PY@tc##1{\textcolor[rgb]{0.41,0.47,0.13}{##1}}}
\@namedef{PY@tok@nt}{\let\PY@bf=\textbf\def\PY@tc##1{\textcolor[rgb]{0.00,0.50,0.00}{##1}}}
\@namedef{PY@tok@nd}{\def\PY@tc##1{\textcolor[rgb]{0.67,0.13,1.00}{##1}}}
\@namedef{PY@tok@s}{\def\PY@tc##1{\textcolor[rgb]{0.73,0.13,0.13}{##1}}}
\@namedef{PY@tok@sd}{\let\PY@it=\textit\def\PY@tc##1{\textcolor[rgb]{0.73,0.13,0.13}{##1}}}
\@namedef{PY@tok@si}{\let\PY@bf=\textbf\def\PY@tc##1{\textcolor[rgb]{0.64,0.35,0.47}{##1}}}
\@namedef{PY@tok@se}{\let\PY@bf=\textbf\def\PY@tc##1{\textcolor[rgb]{0.67,0.36,0.12}{##1}}}
\@namedef{PY@tok@sr}{\def\PY@tc##1{\textcolor[rgb]{0.64,0.35,0.47}{##1}}}
\@namedef{PY@tok@ss}{\def\PY@tc##1{\textcolor[rgb]{0.10,0.09,0.49}{##1}}}
\@namedef{PY@tok@sx}{\def\PY@tc##1{\textcolor[rgb]{0.00,0.50,0.00}{##1}}}
\@namedef{PY@tok@m}{\def\PY@tc##1{\textcolor[rgb]{0.40,0.40,0.40}{##1}}}
\@namedef{PY@tok@gh}{\let\PY@bf=\textbf\def\PY@tc##1{\textcolor[rgb]{0.00,0.00,0.50}{##1}}}
\@namedef{PY@tok@gu}{\let\PY@bf=\textbf\def\PY@tc##1{\textcolor[rgb]{0.50,0.00,0.50}{##1}}}
\@namedef{PY@tok@gd}{\def\PY@tc##1{\textcolor[rgb]{0.63,0.00,0.00}{##1}}}
\@namedef{PY@tok@gi}{\def\PY@tc##1{\textcolor[rgb]{0.00,0.52,0.00}{##1}}}
\@namedef{PY@tok@gr}{\def\PY@tc##1{\textcolor[rgb]{0.89,0.00,0.00}{##1}}}
\@namedef{PY@tok@ge}{\let\PY@it=\textit}
\@namedef{PY@tok@gs}{\let\PY@bf=\textbf}
\@namedef{PY@tok@gp}{\let\PY@bf=\textbf\def\PY@tc##1{\textcolor[rgb]{0.00,0.00,0.50}{##1}}}
\@namedef{PY@tok@go}{\def\PY@tc##1{\textcolor[rgb]{0.44,0.44,0.44}{##1}}}
\@namedef{PY@tok@gt}{\def\PY@tc##1{\textcolor[rgb]{0.00,0.27,0.87}{##1}}}
\@namedef{PY@tok@err}{\def\PY@bc##1{{\setlength{\fboxsep}{\string -\fboxrule}\fcolorbox[rgb]{1.00,0.00,0.00}{1,1,1}{\strut ##1}}}}
\@namedef{PY@tok@kc}{\let\PY@bf=\textbf\def\PY@tc##1{\textcolor[rgb]{0.00,0.50,0.00}{##1}}}
\@namedef{PY@tok@kd}{\let\PY@bf=\textbf\def\PY@tc##1{\textcolor[rgb]{0.00,0.50,0.00}{##1}}}
\@namedef{PY@tok@kn}{\let\PY@bf=\textbf\def\PY@tc##1{\textcolor[rgb]{0.00,0.50,0.00}{##1}}}
\@namedef{PY@tok@kr}{\let\PY@bf=\textbf\def\PY@tc##1{\textcolor[rgb]{0.00,0.50,0.00}{##1}}}
\@namedef{PY@tok@bp}{\def\PY@tc##1{\textcolor[rgb]{0.00,0.50,0.00}{##1}}}
\@namedef{PY@tok@fm}{\def\PY@tc##1{\textcolor[rgb]{0.00,0.00,1.00}{##1}}}
\@namedef{PY@tok@vc}{\def\PY@tc##1{\textcolor[rgb]{0.10,0.09,0.49}{##1}}}
\@namedef{PY@tok@vg}{\def\PY@tc##1{\textcolor[rgb]{0.10,0.09,0.49}{##1}}}
\@namedef{PY@tok@vi}{\def\PY@tc##1{\textcolor[rgb]{0.10,0.09,0.49}{##1}}}
\@namedef{PY@tok@vm}{\def\PY@tc##1{\textcolor[rgb]{0.10,0.09,0.49}{##1}}}
\@namedef{PY@tok@sa}{\def\PY@tc##1{\textcolor[rgb]{0.73,0.13,0.13}{##1}}}
\@namedef{PY@tok@sb}{\def\PY@tc##1{\textcolor[rgb]{0.73,0.13,0.13}{##1}}}
\@namedef{PY@tok@sc}{\def\PY@tc##1{\textcolor[rgb]{0.73,0.13,0.13}{##1}}}
\@namedef{PY@tok@dl}{\def\PY@tc##1{\textcolor[rgb]{0.73,0.13,0.13}{##1}}}
\@namedef{PY@tok@s2}{\def\PY@tc##1{\textcolor[rgb]{0.73,0.13,0.13}{##1}}}
\@namedef{PY@tok@sh}{\def\PY@tc##1{\textcolor[rgb]{0.73,0.13,0.13}{##1}}}
\@namedef{PY@tok@s1}{\def\PY@tc##1{\textcolor[rgb]{0.73,0.13,0.13}{##1}}}
\@namedef{PY@tok@mb}{\def\PY@tc##1{\textcolor[rgb]{0.40,0.40,0.40}{##1}}}
\@namedef{PY@tok@mf}{\def\PY@tc##1{\textcolor[rgb]{0.40,0.40,0.40}{##1}}}
\@namedef{PY@tok@mh}{\def\PY@tc##1{\textcolor[rgb]{0.40,0.40,0.40}{##1}}}
\@namedef{PY@tok@mi}{\def\PY@tc##1{\textcolor[rgb]{0.40,0.40,0.40}{##1}}}
\@namedef{PY@tok@il}{\def\PY@tc##1{\textcolor[rgb]{0.40,0.40,0.40}{##1}}}
\@namedef{PY@tok@mo}{\def\PY@tc##1{\textcolor[rgb]{0.40,0.40,0.40}{##1}}}
\@namedef{PY@tok@ch}{\let\PY@it=\textit\def\PY@tc##1{\textcolor[rgb]{0.24,0.48,0.48}{##1}}}
\@namedef{PY@tok@cm}{\let\PY@it=\textit\def\PY@tc##1{\textcolor[rgb]{0.24,0.48,0.48}{##1}}}
\@namedef{PY@tok@cpf}{\let\PY@it=\textit\def\PY@tc##1{\textcolor[rgb]{0.24,0.48,0.48}{##1}}}
\@namedef{PY@tok@c1}{\let\PY@it=\textit\def\PY@tc##1{\textcolor[rgb]{0.24,0.48,0.48}{##1}}}
\@namedef{PY@tok@cs}{\let\PY@it=\textit\def\PY@tc##1{\textcolor[rgb]{0.24,0.48,0.48}{##1}}}

\def\PYZbs{\char`\\}
\def\PYZus{\char`\_}
\def\PYZob{\char`\{}
\def\PYZcb{\char`\}}
\def\PYZca{\char`\^}
\def\PYZam{\char`\&}
\def\PYZlt{\char`\<}
\def\PYZgt{\char`\>}
\def\PYZsh{\char`\#}
\def\PYZpc{\char`\%}
\def\PYZdl{\char`\$}
\def\PYZhy{\char`\-}
\def\PYZsq{\char`\'}
\def\PYZdq{\char`\"}
\def\PYZti{\char`\~}
% for compatibility with earlier versions
\def\PYZat{@}
\def\PYZlb{[}
\def\PYZrb{]}
\makeatother


% For linebreaks inside Verbatim environment from package fancyvrb. 
\makeatletter
    \newbox\Wrappedcontinuationbox 
    \newbox\Wrappedvisiblespacebox 
    \newcommand*\Wrappedvisiblespace {\textcolor{red}{\textvisiblespace}} 
    \newcommand*\Wrappedcontinuationsymbol {\textcolor{red}{\llap{\tiny$\m@th\hookrightarrow$}}} 
    \newcommand*\Wrappedcontinuationindent {3ex } 
    \newcommand*\Wrappedafterbreak {\kern\Wrappedcontinuationindent\copy\Wrappedcontinuationbox} 
    % Take advantage of the already applied Pygments mark-up to insert 
    % potential linebreaks for TeX processing. 
    %        {, <, #, %, $, ' and ": go to next line. 
    %        _, }, ^, &, >, - and ~: stay at end of broken line. 
    % Use of \textquotesingle for straight quote. 
    \newcommand*\Wrappedbreaksatspecials {% 
        \def\PYGZus{\discretionary{\char`\_}{\Wrappedafterbreak}{\char`\_}}% 
        \def\PYGZob{\discretionary{}{\Wrappedafterbreak\char`\{}{\char`\{}}% 
        \def\PYGZcb{\discretionary{\char`\}}{\Wrappedafterbreak}{\char`\}}}% 
        \def\PYGZca{\discretionary{\char`\^}{\Wrappedafterbreak}{\char`\^}}% 
        \def\PYGZam{\discretionary{\char`\&}{\Wrappedafterbreak}{\char`\&}}% 
        \def\PYGZlt{\discretionary{}{\Wrappedafterbreak\char`\<}{\char`\<}}% 
        \def\PYGZgt{\discretionary{\char`\>}{\Wrappedafterbreak}{\char`\>}}% 
        \def\PYGZsh{\discretionary{}{\Wrappedafterbreak\char`\#}{\char`\#}}% 
        \def\PYGZpc{\discretionary{}{\Wrappedafterbreak\char`\%}{\char`\%}}% 
        \def\PYGZdl{\discretionary{}{\Wrappedafterbreak\char`\$}{\char`\$}}% 
        \def\PYGZhy{\discretionary{\char`\-}{\Wrappedafterbreak}{\char`\-}}% 
        \def\PYGZsq{\discretionary{}{\Wrappedafterbreak\textquotesingle}{\textquotesingle}}% 
        \def\PYGZdq{\discretionary{}{\Wrappedafterbreak\char`\"}{\char`\"}}% 
        \def\PYGZti{\discretionary{\char`\~}{\Wrappedafterbreak}{\char`\~}}% 
    } 
    % Some characters . , ; ? ! / are not pygmentized. 
    % This macro makes them "active" and they will insert potential linebreaks 
    \newcommand*\Wrappedbreaksatpunct {% 
        \lccode`\~`\.\lowercase{\def~}{\discretionary{\hbox{\char`\.}}{\Wrappedafterbreak}{\hbox{\char`\.}}}% 
        \lccode`\~`\,\lowercase{\def~}{\discretionary{\hbox{\char`\,}}{\Wrappedafterbreak}{\hbox{\char`\,}}}% 
        \lccode`\~`\;\lowercase{\def~}{\discretionary{\hbox{\char`\;}}{\Wrappedafterbreak}{\hbox{\char`\;}}}% 
        \lccode`\~`\:\lowercase{\def~}{\discretionary{\hbox{\char`\:}}{\Wrappedafterbreak}{\hbox{\char`\:}}}% 
        \lccode`\~`\?\lowercase{\def~}{\discretionary{\hbox{\char`\?}}{\Wrappedafterbreak}{\hbox{\char`\?}}}% 
        \lccode`\~`\!\lowercase{\def~}{\discretionary{\hbox{\char`\!}}{\Wrappedafterbreak}{\hbox{\char`\!}}}% 
        \lccode`\~`\/\lowercase{\def~}{\discretionary{\hbox{\char`\/}}{\Wrappedafterbreak}{\hbox{\char`\/}}}% 
        \catcode`\.\active
        \catcode`\,\active 
        \catcode`\;\active
        \catcode`\:\active
        \catcode`\?\active
        \catcode`\!\active
        \catcode`\/\active 
        \lccode`\~`\~ 	
    }
\makeatother

\let\OriginalVerbatim=\Verbatim
\makeatletter
\renewcommand{\Verbatim}[1][1]{%
    %\parskip\z@skip
    \sbox\Wrappedcontinuationbox {\Wrappedcontinuationsymbol}%
    \sbox\Wrappedvisiblespacebox {\FV@SetupFont\Wrappedvisiblespace}%
    \def\FancyVerbFormatLine ##1{\hsize\linewidth
        \vtop{\raggedright\hyphenpenalty\z@\exhyphenpenalty\z@
            \doublehyphendemerits\z@\finalhyphendemerits\z@
            \strut ##1\strut}%
    }%
    % If the linebreak is at a space, the latter will be displayed as visible
    % space at end of first line, and a continuation symbol starts next line.
    % Stretch/shrink are however usually zero for typewriter font.
    \def\FV@Space {%
        \nobreak\hskip\z@ plus\fontdimen3\font minus\fontdimen4\font
        \discretionary{\copy\Wrappedvisiblespacebox}{\Wrappedafterbreak}
        {\kern\fontdimen2\font}%
    }%
    
    % Allow breaks at special characters using \PYG... macros.
    \Wrappedbreaksatspecials
    % Breaks at punctuation characters . , ; ? ! and / need catcode=\active 	
    \OriginalVerbatim[#1,codes*=\Wrappedbreaksatpunct]%
}
\makeatother

% Exact colors from NB
\definecolor{incolor}{HTML}{303F9F}
\definecolor{outcolor}{HTML}{D84315}
\definecolor{cellborder}{HTML}{CFCFCF}
\definecolor{cellbackground}{HTML}{F7F7F7}

% prompt
\makeatletter
\newcommand{\boxspacing}{\kern\kvtcb@left@rule\kern\kvtcb@boxsep}
\makeatother
\newcommand{\prompt}[4]{
    {\ttfamily\llap{{\color{#2}[#3]:\hspace{3pt}#4}}\vspace{-\baselineskip}}
}



% Prevent overflowing lines due to hard-to-break entities
\sloppy 
% Setup hyperref package
\hypersetup{
  breaklinks=true,  % so long urls are correctly broken across lines
  colorlinks=true,
  urlcolor=urlcolor,
  linkcolor=linkcolor,
  citecolor=citecolor,
  }
% Slightly bigger margins than the latex defaults

\geometry{verbose,tmargin=1in,bmargin=1in,lmargin=1in,rmargin=1in}














\pagestyle{fancy}
\fancypagestyle{plain}{%
\fancyhf{}
\lhead{\footnotesize\itshape\bfseries\rightmark}
\rhead{\footnotesize\itshape\bfseries\leftmark}
}










% macros
\newcommand{\Q}{\mathbb Q}
\newcommand{\I}{\mathbb I}
\newcommand{\R}{\mathbb R}
\newcommand{\N}{\mathbb N}
\newcommand{\Z}{\mathbb Z}
\newcommand{\C}{\mathbb C}
\newcommand{\K}{\mathbb K}
\newcommand{\M}{\mathcal M}
\newcommand{\MR}{\mathcal M_{nn}(\R)}
\newcommand{\MC}{\mathcal M_{nn}(\C)}

\newcommand{\MK}{\mathcal M_{nn}(\K)}
\newcommand{\ds}{\displaystyle}
\DeclareMathOperator{\tr}{tr}
\DeclareMathOperator{\Ker}{Ker}
\DeclareMathOperator{\Ima}{Im}
\DeclareMathOperator{\rango}{r}
%Teoremas, Lemas, etc.
\theoremstyle{plain}
\newtheorem{teo}{Teorema}
\newtheorem*{teo*}{Teorema}
\newtheorem{lem}{Lema}
\newtheorem{prop}{Proposici\'on}
\newtheorem{cor}{Corolario}

\theoremstyle{definition}
\newtheorem*{defi}{Definici\'on}
\newtheorem{solu}{Solución}

\newenvironment{sol}{\begin{framed}
\begin{solu}
\leavevmode
}{\end{solu}\end{framed}}



\newcommand{\cin}{\operatorname{cint}}
\newcommand{\horrule}[1]{\noindent\rule{\linewidth}{#1}}



\begin{document}
%Encabezado
\fancyhead[L]{ }
\fancyhead[R]{Institute of Science and Technology Austria} 

% \begin{minipage}[t]{10.0 cm}
% \vspace{-6.5ex}
% \textbf{ } \\
% \textbf{} \\
% \textbf{Student: } Mau Rojas 
%    \end{minipage}%
  %  \hfill
  %  \begin{minipage}[t]{3.2 cm}
  %    \vspace{-6.5ex}
  %    \raggedright\includegraphics[scale=1.0]{ISTA_LOGO.png} 
  %  \end{minipage}
\begin{center}
\LARGE\textbf{Hessian terms  \\}
\vspace{3 mm}
\end{center}

    
The idea now is that for every energy term i will add the Hessian terms that will be needed.

I already have the Hessian of the area  i also need to add later the hessian of the Dihedral angle



\begin{equation}
    u = p_2 - p_1 \quad v = p_3 - p_1 \quad w = p_4 - p_1
\end{equation}

The idea now is to write the terms of the Hessian here 


We consider also that 
\begin{align}
    g &= \det(u,v,w)||u|| \\
    h &= \langle u \times w, v \times u \rangle  \\
    f &= \arctan2(g,h) \\
    r &= g^2 + h^2
\end{align}

Now we calculate the derivatives 

\begin{align}
    f_u &= f_h h_u + f_g g_u \\
    &= \cfrac{1}{r}\left( -g u^T ([w]_x[v]_x + [v]_x [w]_x) + h( ||u|| (v\times w)^T + det(u,v,w) \cfrac{u^T}{||u||})  \right)
\end{align}

\begin{align}
    f_v &= f_h h_v + f_g g_v \\
    &= \cfrac{1}{r}\left( -g w^T[u]_x^2 + h||u|| (w\times u)^T\right)
\end{align}
\begin{align}
    f_w &= f_h h_w + f_g g_w \\
    &= \cfrac{1}{r}\left( -g v^T[u]_x^2 + h||u|| (u\times v)^T\right)
\end{align}

Now that we have differentiated with all the variables we can say how is the gradient for every point.


\begin{equation}
    f_{p_1} = -f_u-f_v-f_w \quad f_{p_2} = f_u \quad f_{p_3} = f_v \quad f_{p_4} = f_w
\end{equation}

And those are the terms of the gradient of $f$. Lets remember that f is the dihedral angle.

Now we will do the Hessian of the dihedral angle 


Considering  
\begin{equation}
 l = v\times w \quad B = [w]_x[v]_x + [v]_x[w]_x  \quad \alpha = \det(u,v,w)
\end{equation}

\begin{align}
    f_{p_2p_2}  &= \cfrac{g^2-h^2}{r^2} \left[ ||u|| (Bu \otimes l + l \otimes Bu ) + \cfrac{\alpha}{||u||}\left( Bu \otimes u + u \otimes Bu \right)    \right]\\
    &+\cfrac{2gh}{r^2}\left[ Bu\otimes Bu - ||u||^2 l \otimes l  \right] + \cfrac{1}{r^2}\left(-2gh\alpha + \cfrac{hr}{||u||}\right)\left(l\otimes u + u \otimes l \right)\\
    & - \left(\cfrac{2gh\alpha^2}{r^2 ||u||^2} - \cfrac{h\alpha r }{r^2 ||u||^3} \right)(u \otimes u) - \cfrac{g}{r}B + \cfrac{h\alpha }{||u|| r} \mathcal{I}
\end{align}


This one is the ugliest actually, the other ones are nicer 

First $f_{p_3p_3}$ courtesy of Christian, assuming this 
\begin{equation}
    a = w\times u \quad b = u \times (u\times w) = a\times u 
\end{equation}

\begin{align}
    f_{p_3p_3} = \cfrac{2gh }{r^2}\left(b\otimes b - ||u||^2 a\otimes a \right) + \cfrac{||u||(g^2-h^2)}{r^2} (a\otimes b + b\otimes a)
\end{align}


The next one is $f_{p_4p_4}$.
In this case we consider 
\begin{equation}
    d = u\times v \quad k = u \times (u\times v) = u\times d
\end{equation}

And the terms are 

\begin{align}
    f_{p_4p_4} = \cfrac{2gh}{r^2}\left(k\otimes k - ||u||^2 d\otimes d \right) + \cfrac{||u||(g^2-h^2)}{r^2} (d\otimes k + k\otimes d)
\end{align}

And those are the diagonal terms, now we will do the off diagonals.

We will define some things also 

\begin{equation}
    C = (-2[w]_x[u]_x + [u]_x[w]_x)
\end{equation}

So the next term is $f_{p_2p_3}$

\begin{align}
    f_{p_2 p_3}&= \cfrac{g^2-h^2}{r^2}\left[  -||u||Cv \otimes a + ||u|| l \otimes b  + \cfrac{\alpha }{||u||} u \otimes b  \right]\\
    & +\cfrac{2gh }{r^2}\left[Cv \otimes b + ||u||^2 l\otimes a + \alpha u \otimes a \right] - \cfrac{g}{r}C
\end{align}


And now we have $f_{p_2p_4}$

We define a matrix as  
\begin{equation}
    D = -2[v]_x[u]_x + [u]_x[v]_x
\end{equation}

\begin{align}
    f_{p_2 p_4} &= \cfrac{g^2-h^2}{r^2}\left[||u|| Dw \otimes d + ||i|| l \otimes (u\times d) + \cfrac{\alpha}{||u||}u \otimes (u\times d ) \right]\\
    & +\cfrac{2gh}{r^2}\left[ D w \otimes (u\times d) - ||u||l \otimes d  - \alpha u \otimes d \right] - \cfrac{g}{r}D + \cfrac{h}{r}\left( ||u||[v]_x + \cfrac{1}{||u||} u \otimes d \right)
\end{align}


And now we have $f_{p_4p_3}$

\begin{align}
    f_{p_4p_3}&= \cfrac{g^2-h^2}{r^2}\left[ -\? ||u|| k \otimes a + ||u||d \otimes b   \right]\\
    &+ \cfrac{2gh}{r^2}\left[  ||u||^2 d\otimes a + k \otimes b  \right] - \cfrac{g}{r}[u]_x^2 +\cfrac{h||u||}{r} [u]_x
\end{align}


Now we can do the rest of the terms 

\begin{align}
    f_{p_1p_1 } &= (-f_{p_2} -f_{p_3} -f_{p_4})_p_1\\
                &= f_{p_2p_2} + f_{p_2p_3} + f_{p_2p_4} + f_{p_3p_2} +f_{p_3p_3} + f_{p_3p_4} + f_{p_4p_2} + f_{p_4p_3} + f_{p_4p_4}
\end{align}

Lets do the rest 
 
\begin{align}
    f_{p_1p_2} &= (-f_{p_2} -f_{p_3} -f_{p_4})_p_2\\
                &= -f_{p_2p_2} - f_{p_3p_2} - f_{p_4p_2} 
\end{align}


\begin{align}
    f_{p_1p_3} &= (-f_{p_2} -f_{p_3} -f_{p_4})_p_3\\
                &= -f_{p_2p_3} - f_{p_3p_3} - f_{p_4p_3}
\end{align}

\begin{align}
    f_{p_1p_4} &= (-f_{p_2} -f_{p_3} -f_{p_4})_p_4\\
                &= -f_{p_2p_4} - f_{p_3p_4} - f_{p_4p_4}
\end{align}

\newpage 
 

Ok, now lets write the rest of the hessians. To take a small break i will write the Hessian of an edge 


Let $u = p_2-p_1$


\begin{equation}
    l(u) = \vert| u \vert| 
\end{equation}

From out previous experience we know 
\begin{equation}
    l_{p1} = -l_u \quad l_{p2} = l_u
\end{equation}

Thefore the term of the gradient is 

\begin{equation}
    l_u  = \cfrac{u^T}{||u||} \quad l_{p1} = -\cfrac{u^T}{||u||} \quad l_{p2} = \cfrac{u^T}{||u||}
\end{equation}

And now the for the Hessian we now 
\begin{equation}
    l_{uu} = l_{p_1p_1 } = l+{p_2p_2} = \cfrac{1}{||u||} \left( \mathcal{I} - \cfrac{u\otimes u}{||u||^2} \right)
\end{equation}

Now for the off diagonal terms, hmm... in this case i think they just have a minus sign at the front.

\begin{equation}
    l_{p_1p_2} =  -(l_u)_u = -\cfrac{1}{||u||} \left( \mathcal{I} - \cfrac{u\otimes u}{||u||^2} \right)
\end{equation}



\newpage 

Now we can do the Hessian of the area, we will use the same notation as before, so we have


\begin{equation}
    u = p_2 - p_1 \quad v = p_3 - p_1 \quad u_{p_1} = -id \quad u_{p_2} = id \quad v_{p_1} = -id \quad v_{p_3} = id
\end{equation}

Let $z  = u \times v \quad z_v = [u]_{\times} \quad z_u = -[v]_{\times}$

Considering $A = \vert|z\vert|$ (This is actually twice the area but we just need to remember to divide the results by two at the the end)

\begin{equation}
    A_z = \cfrac{z^T}{\vert| z \vert|} 
\end{equation}


\begin{align}
    A_u - A_z z_u &= -\cfrac{1}{\vert| z\vert|}z^T [v]_\times = \cfrac{1}{\vert| z\vert|}([v]_\times z)^T\\
    &=\cfrac{1}{\vert| z \vert|} (v\times z )^T
\end{align}
\begin{align}
    A_v = A_z z_v & = \cfrac{1}{\vert| z \vert|}z^T [u]_\times = \cfrac{1}{\vert| z \vert|}([u]_\times^T z)^T\\
    &= -\cfrac{1}{\vert| z \vert|} (u\times z)^T
\end{align}


Now that we have the partial derivatives we can write the gradient of the area

\begin{align}
    A_{p_1} &= A_u u_{p_1} + A_v v_{p_1} \\
    &=  -A_u - A_v \\
    &= -\cfrac{1}{\vert| z \vert|} ((u-v)\times z)^T\\
    &= -\cfrac{1}{\vert z \vert|} ((p_2-p_3)\times z)^T 
\end{align}

\begin{align}
    A_{p_2} &= \cfrac{1}{\vert| z |vert}((p_3-p_1)\times z)^T
\end{align}
\begin{align}
    A_{p_3} &= \cfrac{1}{\vert| z \vert|}((p_1-p_2)\times z)^T
\end{align}


Perfect, now that we have the first derivatives its time to write down the second derivatives.


Thanks to Christian we got this 

\begin{equation}
    A_{p_1p_1} = \cfrac{1}{\vert| z  \vert|} \left(  \cfrac{-((u-v)\times z)\cdot ((u-v)\times z)^T }{\vert| z \vert|^2} - [u-v]_\times^2     \right)
\end{equation}

\begin{equation}
    A_{p_2p_2} = -\cfrac{1}{\vert| z \vert|}\left( \cfrac{1}{\vert| z \vect|^2}((p_3-p_1)\times z)\outer((p_3-p_1)\times z) + [p_3-p_1]_\times^2   \right)
\end{equation}

\begin{equation}
    A_{p_3p_3} = -\cfrac{1}{\vert| z \vert|}\left( \cfrac{1}{\vert| z \vect|^2}((p_1-p_2)\times z)\outer((p_1-p_2)\times z) + [p_1-p_2]_\times^2   \right)
\end{equation}

\begin{equation}
    A_{p_2p_3} = \cfrac{1}{\vert| z\vert|}\left([p_3-p_1]_\times (id - \cfrac{z\outer z}{\vert| z \vert|^2})[p_2-p_1]_\times - [z]_\times    \right)
\end{equation}

\begin{equation}
    A_{p_3p_1} = \cfrac{1}{\vert| z \vert|}\left( [p_1-p_2]_\times (id - \cfrac{z\outer z}{\vert| z \vert|^2})[p_3-p_2]_\times - [z]_\times   \right)
\end{equation}

\begin{equation}
    A_{p_1p_2} = \cfrac{1}{\vert| z \vert|}\left( [p_2-p_3]_\times (id - \cfrac{z\outer z}{\vert| z \vert|^2})[p_1-p_3]_\times - [z]_\times   \right)
\end{equation}






I think those are all the terms that i need for now. I still need to add the interactions with the Bead but lets do that later because we require considering dot products :( 





\chapter{Assembling the Hessian}

Ok so this is the part where we dialog the structure of the code to get my Hessian. 

According to Christian I should create functions that take a vector and return the Block hessians.

For the area hessian i should have 

\begin{itemize}
    \item \texttt{gradient_area([p1 p2, p3])} This function takes a 9-dimensional vector with the coordinates of the three vertices and returns a  9-vector that contains the gradient of the area with respect to each of the vertices.
    \item \texttt{hessian_area([p1 p2, p3])} This function takes a 9-dimensional vector with the coordinates of the three vertices and returns a 9x9 matrix that contains the Hessian of the area with respect to each of the vertices.
\end{itemize}

This functions should be implemented simply because there will be a bigger function that will ensamble the rest of the info.



\end{document}